\needspace{10\baselineskip}
\SourceLeaf{316}
\begin{center}
\pdfbookmark[1]{14. Verbesserung der Elemente durch vier Oppositionen}{zusatz-014}
[14.]
\end{center}
\begin{center}
\{Über die Verbesserung der Elemente eines Planeten\\
durch vier beobachtete Oppositionen.\}
\end{center}
Wenn die Elemente eines Planeten schon so genau bekannt sind, dass die
angenommenen Werthe von den wahren nur sehr wenig abweichen und man hat die
Beobachtungen von vier Oppositionen, so lässt sich die Verbesserung der
Elemente leichter und einfacher vollbringen, als wenn die vier gegebenen
Örter des Planeten beliebig genommen wären. Das Theoretische dieser
Untersuchung ist schon in der Abhandlung »Disquisitio de elementis ellipticis
Palladis« mitgetheilt; hier folgen einige praktische Bemerkungen, die überhaupt
bei der Berechnung der Oppositionen und dergl. brauchbar sind.

Zuerst muss die Epoche verbessert werden. Man berechnet also aus einer
beliebigen Beobachtung, am besten aus einer solchen, die zu einer der mittlern
Oppositionen gehört, die geocentrische Länge und Breite des Planeten, aus
diesen mit Hülfe des aus der unveränderten Epoche durch die bekannten Formeln
erhaltenen Radius Vector, die heliocentrische Länge und Breite, aus diesen die
Länge in der Bahn. Nachdem an diese Länge die Störungen (wenn schon Formeln
für sie vorhanden sind) angebracht, leitet man aus ihr die mittlere Anomalie
her, welche mit der aus der unveränderten Epoche erhaltenen verglichen, die
Änderung der Epoche gibt. Sind keine Beobachtungen, sondern nur die
berechneten Längen und Breiten zu der Zeit der vier Oppositionen gegeben, so
wird man eine dieser Längen zu wählen haben; besser ist es indessen, die
Beobachtungen selbst vorzunehmen, als sich auf die Berechnungen Anderer zu
verlassen.

Mit der auf solche Art verbesserten Epoche und den übrigen unverändert
gelassenen Elementen werden nun die vier Oppositionen vorläufig berechnet,
wobei man sich mit hinlänglicher Genauigkeit früherer Berechnungen bedienen
kann. Um eine Opposition am bequemsten berechnen zu können, verfährt man
folgendermaassen.

Man wählt, wenn von mehrern Orten Beobachtungen da sind, die Beobachtungen
desjenigen Orts, der die meisten geliefert hat, bestimmt aus den angegebenen
Zeiten des Durchgangs durch den Mittagskreis, die Culminationszeiten des
Planeten von zwei zu zwei Tagen für die ganze Zeit, welche alle Beobachtungen
umfassen. Durch eine vorläufige Rechnung, für einen in der Mitte der
Beobachtungszeit liegenden Tag, bestimmt man die Entfernung des Planeten von
der Erde und aus dieser die Zeit, welche das Licht gebraucht, um von dem
Planeten auf die Erde zu kommen, welche Zeit von den Culminationszeiten
abzuziehen ist; hiedurch hat man für die Aberration Rechnung gehalten. Aus
der verbesserten Epoche leitet man nun durch Abziehung der für die Zeit der
Epoche stattfindenden Länge des Perihels die mittlere Anomalie her und
bestimmt aus dieser mittelst der täglichen tropischen Bewegung die mittlern
Anomalien für die angenommenen Tage und aus diesen die Rectascensionen und
Declinationen. Um für die Nutation Rechnung zu halten, wendet man nicht die
mittlere durch die Elemente gegebene Länge des Knotens, sondern die wahre an,
die man durch Hinzufügung der Nutation in der Länge (\textsc{Bohnenbergers}
Astron. p. 672) zu jener erhält; die Parallaxe aber bringt man nach den in der
Theoria motus Corpor. Coel. art. 59 oder 70 erklärten Methode an. Übrigens
sieht man leicht ein, weshalb die Rechnung von zwei zu zwei Tagen
durchgeführt ist, auch wenn an einigen derselben nicht sollte beobachtet sein;
man kann nemlich die Rechnung sehr bequem durch Differenzen prüfen, so dass
es nicht möglich ist, einen Fehler zu übersehen, und dann vermittelst einer
scharfen Interpolation die Berechnungen für \textls[150]{alle} Tage erhalten.

Die auf solche Weise erhaltenen Rectascensionen und Declinationen vergleicht
man mit den beobachteten, bestimmt hieraus die mittlere Abweichung sowohl der
Rectascension als der Declination und indem man diese zu den berechneten
hinzufügt, leitet man für den Tag der Opposition, welcher sehr bald
\SourceLeaf{317}durch Vergleichung der Beobachtungen mit den berechneten Sonnenlängen gefunden wird, die geocentrische Länge her, bestimmt eben dieselbe für einen um eine beliebige Grösse (etwa einige Stunden) von der vorhin angenommenen Zeit abweichenden Zeitpunkt und sucht nun durch Interpolation den Augenblick zu bestimmen, in welchem die Länge der Sonne und die des Planeten genau um \(180^{0}\) von einander verschieden sind. Bei der Berechnung der Sonne ist übrigens für die durch die Aberration des Planeten verbesserte Culminationszeit die \textls[150]{wahre} Länge zu suchen, also die durch die Tafeln gefundene, welche die Aberration der Sonne bereits einschliesst, um diese zu vermehren.

Sind nur wenige Beobachtungen vorhanden (etwa bis sechs), so kann man auch, wenn man es für vortheilhafter hält, die Rectascensionen und Declinationen bloss für die Beobachtungstage berechnen.

Bei der vorläufigen Berechnung der Oppositionen kann man allenfalls die Aberration, Nutation und Parallaxe vernachlässigen; will man aber die Oppositionen nur einmal berechnen, so muss nothwendig auf sie Rücksicht genommen werden.

Aus den gefundenen Längen und Breiten zu der Zeit der vier Oppositionen leitet man nun nach der in der Abhandlung de Elementis Palladis angeführten Methode die Verbesserungen der einzelnen Elemente her.

\par\medskip\begin{center}\rule{.22\textwidth}{.3pt}\end{center}
